%%%%%%%%%%%%%%%%%疑似コードを書くためのパッケージと設定%%%%%%%%%%%%%%%%%%%%%
%listingsパッケージの設定
\lstset{%
  language={Pascal},
  basicstyle={\small},%
  identifierstyle={\small},%
  commentstyle={\small\itshape\color[rgb]{0,0.5,0}},%
  keywordstyle={\small\bfseries\color[rgb]{0,0,1}},%
  ndkeywordstyle={\small},%
  stringstyle={\small\ttfamily\color[rgb]{1,0,1}},
  frame={tb},
  breaklines=true,
  caption= Algorithm ,
  columns=[l]{fullflexible},%
  numbers=left,%
  xrightmargin=0zw,%
  xleftmargin=3zw,%
  numberstyle={\scriptsize},%
  stepnumber=1,
  numbersep=1zw,%
  lineskip=-0.5ex,%
 escapechar=\@
}
\renewcommand{\lstlistingname}{Algorithm}
%%%%%%%%%%%%%%%%%%%%%%%%%%%%%%%%%%%%%%%%%%%%%%%%%%%%%%%%%%%%%%%%%%%%%





%%%%%%subcaptionの使用によるcaptionの書式の変更を防止：書き換えないこと%%%%%%%
\makeatletter %「@」をアルファベットとして認識する
\let\MYcaption\@makecaption
\makeatother %「@」のアルファベットとして認識終了

\usepackage{subcaption} %subfigの更新版
\captionsetup{compatibility=false} %必要に応じて

\makeatletter
\let\@makecaption\MYcaption
\makeatother
%%%%%%%%%%%%%%%%%%%%%%%%%%%%%%%%%%%%%%%%%%%%%%%%%%%%%%%%%%%%%%%%%%%%%





%********************** <Various setting> **********************%
\newcommand{\argmax}{\mathop{\rm argmax}\limits}
\newcommand{\argmin}{\mathop{\rm argmin}\limits}
\newcommand{\ctext}[1]{\raise0.2ex\hbox{\textcircled{\scriptsize{#1}}}}

\renewcommand{\baselinestretch}{1.1} %行間調節→{1}を調節する
\renewcommand{\bibname}{参考文献} %参考文献のタイトル変更

\makeatletter
\def\Hline{
\noalign{\ifnum0=`}\fi\hrule \@height 1pt \futurelet
\reserved@a\@xhline}
\makeatother
%************************************************************%






%******************** <Setting about figures> ********************%
\makeatletter
\def\fnum@figure{Fig.\thefigure} %図のキャプション文字変更
\def\fnum@table{Table \thetable} %表のキャプション文字変更
\makeatother

\setlength\abovecaptionskip{0pt}  %｛｝で表，図とcaptionの間の間隔調整
%\setlength\belowcaptionskip{0pt}  %{}でcaptionの下の間隔調整

%\renewcommand{\subfigcapskip}{0pt} %キャプションと図の間の間隔調整
%\renewcommand{\subfigbottomskip}{-2pt}

\providecommand{\fig}[5][H] %画像挿入省略のためのマクロ定義
{
\begin{figure}[#1]
\centering
\includegraphics[#2]{#3}
\caption{#4}
\label{#5}
\end{figure}
}
%************************************************************%





%*********************** <Page setting> ***********************%
\oddsidemargin=10mm %奇数ページの左側の余白
\evensidemargin=1mm %偶数ページの左側の余白
\topmargin=-5mm %上部の余白
\textheight=252mm %本文の高さ
\headheight=5mm %ヘッダ記述部分の高さ
\headsep=5mm %ヘッダと本文の間の幅
%\footskip=10mm %本文からフッタの下部までの幅
%************************************************************%





%****************** <Customizing chapter style> ******************%
%\newcommand\chapformat[1]{%
% \parbox[b]{.5\textwidth}{\filleft\bfseries #1}%
% \quad\rule[-12pt]{2pt}{70pt}\quad
% {\fontsize{60}{60}\selectfont\thesection}}
%\titleformat{\chapter}[block]
% {\filleft\normalfont\sffamily}{}{0pt}{\chapformat}

\setcounter{topnumber}{5}%    ページ上部の図表は 5 個まで
\def\topfraction{1.00}%       ページの上 1.00 まで図表で占めて可
\setcounter{bottomnumber}{5}% ページ下部の図表は 5 個まで
\def\bottomfraction{1.00}%    ページの下 1.00 まで図表で占めて可
\setcounter{totalnumber}{10}% ページあたりの図表は 10 個まで
\def\textfraction{0.04}%      ページうち本文が占める割合の下限

\makeatletter
\newcommand{\figcaption}[1]{\def\@captype{figure}\caption{#1}} %figure環境内で\captionコマンドと同じ効果：恐らくtablarとの判別用
\newcommand{\tblcaption}[1]{\def\@captype{table}\caption{#1}} %tblcaption環境内で\captionコマンドと同じ効果：恐らくfigureとの判別用
\makeatother
%************************************************************%





%*************** <Customizing style of itemization> ***************%
\newenvironment{itemize2} %箇条書きの書式をカスタマイズ
{
   \begin{list}{$\bullet$\ \ } %見出し記号/直後の空白を調節
   {
      \setlength{\itemindent}{0pt}
      \setlength{\leftmargin}{3zw} %左のインデント
      \setlength{\rightmargin}{0zw} %右のインデント
      \setlength{\labelsep}{0zw} %黒丸と説明文の間
      \setlength{\labelwidth}{0.5zw} %ラベルの幅
      \setlength{\itemsep}{.35em} %項目ごとの改行幅
      \setlength{\parsep}{0em} %段落での改行幅
      \setlength{\listparindent}{0zw} %段落での一字下り
   }
}
{
   \end{list}
}

\newenvironment{itemize3} %箇条書きの書式をカスタマイズ
{
   \begin{list}{$\bullet$\ \ } %見出し記号/直後の空白を調節
   {
      \setlength{\itemindent}{0pt}
      \setlength{\leftmargin}{3zw} %左のインデント
      \setlength{\rightmargin}{0zw} %右のインデント
      \setlength{\labelsep}{0zw} %黒丸と説明文の間
      \setlength{\labelwidth}{0.5zw} %ラベルの幅
      \setlength{\itemsep}{.4em} %項目ごとの改行幅
      \setlength{\parsep}{0em} %段落での改行幅
      \setlength{\listparindent}{0zw} %段落での一字下り
   }
}
{
   \end{list}
}

\newcounter{enum2} %番号付き箇条書きのカスタマイズ
\newenvironment{enumerate2}{
   \begin{list}
   {
      \arabic{enum2})\ \, %見出し記号/直後の空白を調節
   }
   {%評価値
      \usecounter{enum2}
      \setlength{\itemindent}{0em} %ここは 0 に固定
      \setlength{\leftmargin}{1.5em} %左のインデント
      \setlength{\rightmargin}{0em} %右のインデント
      \setlength{\labelsep}{0em} %黒丸と説明文の間
      \setlength{\labelwidth}{0.5em} %ラベルの幅
      \setlength{\itemsep}{.35em} %項目ごとの改行幅
      \setlength{\parsep}{0em} %段落での改行幅
      \setlength{\listparindent}{0em} %段落での一字下り
   }
}
{
   \end{list}
}

\newcounter{enum3} %番号付き箇条書きのカスタマイズ
\newenvironment{enumerate3}{
   \begin{list}
   {
      (\arabic{enum3})\ \, %見出し記号/直後の空白を調節
   }
   {
      \usecounter{enum3}
      \setlength{\itemindent}{0em} %ここは 0 に固定
      \setlength{\leftmargin}{1.5em} %左のインデント
      \setlength{\rightmargin}{0em} %右のインデント
      \setlength{\labelsep}{0em} %黒丸と説明文の間
      \setlength{\labelwidth}{0.5em} %ラベルの幅
      \setlength{\itemsep}{.5em} %項目ごとの改行幅
      \setlength{\parsep}{0em} %段落での改行幅
      \setlength{\listparindent}{0em} %段落での一字下り
   }
}
{
   \end{list}
}
%************************************************************%